\documentclass[11pt]{article}
\usepackage{times}
\usepackage{fullpage}
\usepackage{url}
\usepackage{hyperref}
\usepackage{fancyhdr}
\usepackage{graphicx}
\usepackage{enumitem}
\usepackage{subcaption}
\usepackage{amsmath, amsfonts, amsthm, fouriernc}
\usepackage{IEEEtrantools}
\usepackage{parskip}

\pagestyle{fancy}
\addtolength{\headheight}{2em}
\addtolength{\headsep}{1.5em}
\lhead{\large{ME 2801}}

\usepackage{sectsty}
\sectionfont{\fontsize{12}{8}\selectfont}

\begin{document}

\begin{center}
  \Large{\bf{Lab Assignment: USV Navigation}}
\end{center}

This assignment accompanies the autonomous navigation experiments on the USV.
You will analyze waypoint-following performance, connect outer-loop navigation
to inner-loop speed and heading control, and reflect on the end-to-end
system behavior.

\textbf{Deliverable.}  A brief lab report (PDF) with figures and MATLAB code
appended.

\bigskip\hrule\bigskip

% ---------------------------------------------------------------
\section*{1\quad Pre-lab: Navigation Architecture}

\begin{enumerate}

\item Describe the two-loop control architecture used for waypoint navigation:
  identify the outer loop (guidance) and inner loops (speed and heading
  control).  Draw a block diagram showing how the loops interact.

\item For the heading control loop, what is the commanded variable?  How does
  the guidance law translate a waypoint bearing error into a heading command?

\item Based on your heading controller gains from the previous lab, predict
  the closed-loop bandwidth of the heading loop.  Explain why bandwidth
  matters for navigation performance.

\end{enumerate}

% ---------------------------------------------------------------
\section*{2\quad Experiment}

\begin{enumerate}[resume]

\item Configure ArduPilot for an autonomous waypoint mission with at least four
  waypoints forming a closed course.  Record the parameter file and the mission
  file.

\item Execute the mission and collect ArduPilot \texttt{.BIN} logs.
  If possible, run at least two trials with different speed commands.

\end{enumerate}

% ---------------------------------------------------------------
\section*{3\quad Post-lab: Data Analysis}

\begin{enumerate}[resume]

\item Process the \texttt{.BIN} log.  Plot the GPS track (latitude/longitude
  converted to local East-North coordinates) and overlay the commanded
  waypoints.

\item Compute the cross-track error (perpendicular distance from the planned
  path) as a function of time.  Plot and report the RMS cross-track error for
  each leg.

\item Plot heading and speed as functions of time.  Identify the time intervals
  corresponding to turns and straight-line segments.

\end{enumerate}

% ---------------------------------------------------------------
\section*{4\quad Discussion}

\begin{enumerate}[resume]

\item How does cross-track error vary with vehicle speed?  Is this consistent
  with your expectation based on the heading-control bandwidth?

\item Identify one factor that limited navigation accuracy in your experiment
  and propose a specific change (to the controller or the guidance law) that
  would reduce it.

\item Reflect on the full-stack design process from this course: plant
  identification $\to$ controller design $\to$ closed-loop analysis $\to$
  field implementation.  What was the most significant gap between theory and
  experiment?

\end{enumerate}

\end{document}
